\documentclass[11pt,a4paper]{article}
\usepackage[margin=24mm]{geometry}
\usepackage{amsmath,amssymb,amsthm,hyperref,booktabs}
\newtheorem{theorem}{Theorem}\newtheorem{lemma}[theorem]{Lemma}\newtheorem{proposition}[theorem]{Proposition}
\newcommand{\gen}[1]{\langle #1\rangle}
\title{The complete exclusion list for simple string C-groups}
\author{Human Compression project: recovery of L011}\date{}
\begin{document}\maketitle
We recover the principal classification of Leemans and Vandenschrick~\cite{LV}.
The complete companion \emph{Projective involutions and support spaces in small
dimension} is an acquired common proof. The finite exclusions below are
reproducible exhaustive computations, with a proved enumeration boundary.
The positive direction retains the deep generation classification attributed
by the source to Nuzhin and Mazurov; it is not reproved by these computations.

\section{The retained positive classification and the intersection property}
We retain the following exact statement of the established rank-three
generation results, as quoted in the introduction of~\cite{LV}: every
nonabelian finite simple group outside
\[
\begin{gathered}
 \operatorname{PSL}_3(q),\quad \operatorname{PSU}_3(q),\quad
 \operatorname{PSL}_4(2^n),\quad \operatorname{PSU}_4(2^n),\\
 \operatorname{PSU}_4(3),\quad\operatorname{PSU}_5(2),\quad A_6,\quad A_7,
 \quad M_{11},\quad M_{22},\quad M_{23},\quad \mathrm{McL}
\end{gathered}\tag{1}
\]
is generated by three involutions with two commuting. Here $q$ ranges over
prime powers and $n\ge1$, and only the nonabelian simple members of the
listed families are relevant. The source attributes the classical and
alternating cases to Nuzhin, the sporadic cases to Mazurov, and the additional
exceptions $\operatorname{PSU}_4(3),\operatorname{PSU}_5(2)$ to Jones and
Ma\v caj (personal communication). Our finite checks independently establish
their negative direction. The positive classification remains borrowed,
including its classification-theoretic hypotheses and original proof costs.
Names in (1) specify abstract isomorphism types; coincidences of different
classical names do not create new cases.

The following elementary proof supplies the rank-three intersection step
in the simple case. The source cites the more general Conder--Oliveros theorem.

\begin{lemma}\label{rank3}
Let a nonabelian finite simple group $G$ be generated by involutions $a,b,c$
with $a,c$ commuting. Then $(a,b,c)$ is a rank-three string C-group marking.
\end{lemma}
\begin{proof}
The generators are distinct and neither adjacent pair commutes: otherwise
the generated group would be dihedral or would have a nontrivial commuting
normal factor. Put $A=\gen{a,b}$, $B=\gen{b,c}$ and $H=A\cap B$.
Both $A,B$ are proper dihedral subgroups. If $H>\gen b$, then $H$ is a
dihedral group, since a subgroup of a dihedral group containing its reflection
$b$ consists of a cyclic rotation subgroup and its coset by $b$.
Let $X=H\cap\gen{ab}$ and $Y=H\cap\gen{bc}$.
If $H$ is not a four-group, its cyclic index-two rotation subgroup is unique,
so $X=Y\ne1$. This subgroup is normal in both $A$ and $B$, since every
subgroup of a dihedral rotation subgroup is normal. It is then a nontrivial
cyclic normal subgroup of $G$, a contradiction.

If $H$ is a four-group and $X=Y$, the same argument applies. Otherwise
write $X=\gen u$, $Y=\gen v$; its third involution is $b=uv$.
The element $u$ is central in $A$ and $v$ is central in $B$. Thus every
one of $a,v$ commutes with every one of $c,u$. It follows that
$G=\gen{a,v}\gen{c,u}$ is a product of two commuting nontrivial normal
subgroups, impossible for a nonabelian simple group.
Consequently $A\cap B=\gen b$.

The endpoint subgroup is $\gen{a,c}\cong C_2^2$. Its intersection with
$A$ contains $a$; if larger, it contains $c$ and would give $G=A$.
The dual assertion holds for $B$. These three intersections, and independence
(which follows since no two involutions generate $G$), give all subset
intersections in rank three.
\end{proof}

\section{A complete finite elimination method}
For a fixed faithful finite permutation group $G$, we enumerate ordered
prefixes $p=(t_0,\ldots,t_{r-1})$ of connected string C-group markings,
up to simultaneous conjugation by $G$. A prefix is retained only if it has
the full intersection property. Define
\[
 H(p)=\gen{t_0,\ldots,t_{r-1}},\quad
 Z(p)=C_G(t_0,\ldots,t_{r-1}),\quad
 E(p)=\{x\in C_G(t_0,\ldots,t_{r-2}):|x|=2\}.
\]
For length one the last set means all involutions of $G$.
The algorithm starts with one representative of each involution conjugacy
class. At $p$ it filters $E(p)$ by $x\notin H(p)$ and
$|t_{r-1}x|>2$, then takes one representative of every $Z(p)$-orbit
under conjugation. For each representative it tests independence and all
subset intersections of the extended prefix. A cheap adjacent-dihedral
intersection check is applied first; it is only an early necessary rejection.
If the extension generates $G$ it is recorded. Otherwise we recurse.
The next eligible set is computed by filtering $E(p)$ for commutation with
$t_{r-1}$, so is the entire involution set of $Z(p)$.

\begin{proposition}\label{complete}
This algorithm terminates and finds every connected string C-group marking
of $G$, up to simultaneous inner conjugation. For a nonabelian simple group,
absence of such a marking proves absence of every string C-group marking.
\end{proposition}
\begin{proof}
Conjugate the first generator to its chosen class representative. Inductively,
a conjugation preserving the already fixed ordered prefix is exactly an
element of $Z(p)$. String commutation places the next generator in $E(p)$;
independence excludes $H(p)$, and connectedness requires its adjacent product
order to exceed two. The candidate set is invariant under $Z(p)$, so selecting
one representative per orbit loses no possible continuation. The independence
and intersection tests reject only prefixes that cannot occur in a C-group.
At every recursion the generated subgroup strictly increases. Its order at
least doubles, so no branch has unbounded length. This proves termination
and exhaustive coverage, without imposing a conjectured rank cutoff.
Finally, a commuting adjacent pair in a C-group splits it as a direct product
of two nontrivial subgroups, by string commutation and the intersection
property. Simplicity excludes this possibility.
\end{proof}

The full intersection test literally constructs $G_I=\gen{t_i:i\in I}$
for every subset and checks $G_I\cap G_J=G_{I\cap J}$ for every pair.
Thus no maximal-parabolic sufficiency criterion or subgroup-order inference
is silently substituted. Empty subgroups are constructed explicitly.
At each node the program also checks that the conjugacy orbits cover the
whole candidate set disjointly.

\section{Exact models and results}
Alternating groups use their natural permutation models; the three Mathieu
groups use GAP's named standard permutation constructors. The McLaughlin
group is the named primitive catalogue group $\operatorname{PrimitiveGroup}(275,1)$.
Their identification with these named groups is retained from the standard
library. Fresh groups made only from copies of the actual permutation generators have
no cached size, stabilizer chain or ambient parent. Verified stabilizer chains
on these fresh groups, with no supplied size or upper bound, check their exact orders.
The library versions and full input provenance are bound separately; order
agreement alone is not asserted to identify an arbitrary permutation group.

For $\operatorname{PSU}_4(3)$ and $\operatorname{PSU}_5(2)$, construct the
standard $\operatorname{SU}_n(q)$ generators over $\mathbb F_{q^2}$ and
obtain their invariant Hermitian form. The program checks its Hermitian
symmetry and nonsingularity, each generator's determinant one and isometry
equation, and constructs \emph{all} isotropic projective points. A canonical
representative has first nonzero coordinate one. There are respectively
280 and 165 points. Fresh uncached copies of their permutation actions have verified orders respectively
3,265,920 and 13,685,760; its kernel has order $\gcd(n,q+1)$.
These equal the classical projective-unitary order formulas. Together with
the containment established by the isometry equations, the order equality
identifies the full projective unitary image. The standard scalar-center
and order formulas are retained classical inputs.

The following table reports the whole prefix counts, by length, in the
completed search. These are counts of conjugacy representatives of proper
prefixes, not counts of polytopes or abstract subgroup isomorphism classes.
No longer prefix exists in these runs; this is an output of exhaustion.
\begin{center}
\begin{tabular}{lrrl}
\toprule
Group&degree&order&proper-prefix counts\\\midrule
$A_6$&6&360&$1,6,18$\\
$A_7$&7&2,520&$1,6,28,1$\\
$M_{11}$&11&7,920&$1,6,32,3$\\
$M_{22}$&22&443,520&$1,8,178,40$\\
$M_{23}$&23&10,200,960&$1,6,166,77,4$\\
$\mathrm{McL}$&275&898,128,000&$1,4,132,65,2$\\
$\operatorname{PSU}_4(3)$&280&3,265,920&$1,9,185,36,2$\\
$\operatorname{PSU}_5(2)$&165&13,685,760&$2,13,145,26,2$\\\bottomrule
\end{tabular}
\end{center}
Every row has zero generating C-group markings. The positive $A_5$ control
has proper-prefix counts $1,3$ and yields six generating rank-three markings
up to inner conjugacy. These finite results, the exhaustive algorithm, and
its group models are required computational premises of the eight-case
exclusion. They are not optional validation of a universal written theorem.

The implementation uses GAP permutation-group computation. Its guaranteed
verification option is
\[
 \texttt{DefaultStabChainOptions.random=1000}.
\]
The marking enumeration is exhaustive; internal stabilizer-chain
construction may use random choices. No claim of random-free group-library
internals is made. The proof assumes the correctness of the explicitly named
GAP operations and group library. No full Lean verification is claimed.

\section{The complete principal contract}
\begin{theorem}\label{classification}
A nonabelian finite simple group admits a string C-group representation if
and only if it is not one of the isomorphism types in (1).
\end{theorem}
\begin{proof}
The retained generation classification and Lemma~\ref{rank3} give a
rank-three representation for every group outside (1). For the two
projective dimension-three families and both even-characteristic
dimension-four families, use the full geometric companion. The eight
remaining cases are excluded by Proposition~\ref{complete} and the complete
finite certificate. These alternatives exhaust exactly the source list.
\end{proof}

Theorem~\ref{classification} is a recovery of the source's principal result,
not a new classification or a priority claim. We also recover its projective
form theorem with an explicit characteristic-two proof and a direct rank-four
form argument. The source's missing finite implementation has been replaced
by the whole auditable computation above; a bare assertion that a computer
can check the groups is not retained as a theorem. Its positive classification
and the classical structural inputs remain attributed and acquired externally.
The previous odd-characteristic four-dimensional phase-nine and phase-eighteen
notes are contextual prior work, not a proof of these characteristic-two cases;
no circular appeal is made to the L024 application, which itself cites L011.
There is no new question solution, numerical factor, or complete formalization.

\begin{thebibliography}{9}
\bibitem{LV} D. Leemans and A. Vandenschrick, \emph{On simple groups that admit
a string C-group representation}, fixed six-page author manuscript,
\url{https://dipot.ulb.ac.be/dspace/bitstream/2013/397490/3/simpleGroups6.pdf}.
Its introduction supplies the exact Nuzhin--Mazurov attribution and the
Jones--Ma\v caj correction; references [5]--[9] identify the retained papers.
\bibitem{CO} M. Conder and D. Oliveros, \emph{The intersection condition for
regular polytopes}, J. Combin. Theory A \textbf{120} (2013), 1291--1304.
\bibitem{GAP} The GAP Group, \emph{GAP}, version 4.15.1, permutation-group
manual, especially stabilizer-chain verification,
\url{https://gap-system.github.io/gap/doc/ref/chap43.html}.
\end{thebibliography}
\end{document}
